\section{工具使用训练：从语义能力到操作能力}

SFT 和 RL 的机制只有落到外部环境，才真正变成 Agent 训练。本节因此把“会生成函数调用文本”与“能在权限、异常和副作用约束下完成任务”分开讨论，并用 synthetic environment 说明轨迹数据怎样被生成、验证和回放。

\subsection{为什么工具使用是 Agent 训练的分水岭}
课程在后半段详细讨论 tool use，原因在于工具调用让模型进入 ``动作有副作用'' 的世界。模型不仅要决定 ``说什么''，还要决定 ``做什么''，并承担执行后果。

\begin{importantbox}{工具调用能力的本质}
\textbf{工具调用不是函数格式学习，而是决策学习。} 需要同时解决何时调用、调用哪个工具、参数如何构造、失败后如何恢复四类问题。
\end{importantbox}

\subsection{合成工具数据与三模型模拟}
上一小节定义了工具决策的四个问题，这里进一步回答数据从哪里来。真实 API 交互昂贵、慢且可能产生不可逆副作用，因此课程讨论用模拟用户、模拟工具和待训练策略构成受控闭环；关键是让模拟器暴露真实错误类型，并让 verifier 检查终局结果，而不只是 JSON 是否可解析。
讲者介绍了通过 synthetic tools 扩展训练空间的做法：用一个代理模型扮演用户、一个模型扮演工具环境、一个模型作为待训练 agent，在闭环中生成大量多样交互轨迹。

\begin{knowledgebox}{三角色模拟的收益}
\begin{itemize}
    \item 扩展工具覆盖面，降低真实工具接入成本；
    \item 更快制造 corner case，提升鲁棒性；
    \item 在可控环境中复现实验，便于 ablation。
\end{itemize}
\end{knowledgebox}

\subsection{合成数据的边界}
合成数据虽高效，但容易产生分布偏移。若工具模拟器过于理想化，模型会学到 ``仿真世界最优策略''，而非真实 API 世界的稳健策略。

\begin{warningbox}{合成工具训练的三类偏差}
\begin{enumerate}
    \item 错误类型偏差：仿真器未覆盖真实系统的脏数据与异常码。
    \item 延迟偏差：真实网络波动和队列抖动在仿真中被简化。
    \item 权限偏差：真实鉴权、配额、审计流程被忽略。
\end{enumerate}
\end{warningbox}

\subsection{本章小结}
工具使用训练决定 Agent 能否在真实环境落地。合成工具是加速器，但必须与真实工具回放结合，否则泛化能力会被高估。


